\documentclass{article}

\usepackage[T1]{fontenc}

\usepackage{booktabs,graphicx,siunitx}

\begin{document}

\section{Factorial ANOVA}

\begin{tabular}{llll}
\toprule
Factor A $\times$ B & N & Mean & Std. Dev. \\
\midrule
Blended $\times$ Female & 40 & 66.63 & 6.73 \\
Flipped $\times$ Female & 40 & 70.88 & 6.55 \\
Lecture $\times$ Female & 40 & 59.58 & 7.45 \\
Blended $\times$ Male & 40 & 64.36 & 6.86 \\
Flipped $\times$ Male & 40 & 70.62 & 6.69 \\
Lecture $\times$ Male & 40 & 61.34 & 7.67 \\
\bottomrule
\end{tabular}

\begin{tabular}{lllllll}
\toprule
Source & SS & df & MS & F & p & partial $\eta$\textsuperscript{2} [95\% CI] \\
\midrule
teaching\_method & 4234.40 & 2 & 2117.20 & 43.14 & $<$.001 & 0.27 [0.19, 1.00] \\
gender & 3.85 & 1 & 3.85 & 0.08 & .780 & 0.00 [0.00, 1.00] \\
teaching\_method $\times$ gender & 162.25 & 2 & 81.12 & 1.65 & .194 & 0.01 [0.00, 1.00] \\
\bottomrule
\end{tabular}

\begin{tabular}{lllll}
\toprule
Contrast & Mdiff & SE & padj & 95\% CI \\
\midrule
Blended - Flipped & $-$5.26 & 1.11 & $<$.001 & [$-$7.87, $-$2.64] \\
Blended - Lecture & 5.03 & 1.11 & $<$.001 & [2.42, 7.64] \\
Flipped - Lecture & 10.29 & 1.11 & $<$.001 & [7.68, 12.90] \\
\bottomrule
\end{tabular}

\includegraphics{figures/01_factorial-anova/figure_interaction.png}

Each factor's F/p is its main effect; the A$\times$B row tests whether the effect of one factor depends on the other. A significant interaction usually takes priority --- read it from the interaction plot before the main effects. A significant effect for any factor with 3+ levels tells you the levels differ somewhere, not which ones --- use the post-hoc / simple-effects (emmeans) table to find the specific pairs. Your significance threshold ($\alpha$) is 0.05.

A two-way ANOVA gave A$\times$B interaction F(2,234)=1.65, p = .194, partial $\eta$\textsuperscript{2}=.01 [.00, 1.00].

Statistical basis: Factorial ANOVA (main effects + interaction) (Fisher, R. A. (1925). Statistical Methods for Research Workers. Oliver \& Boyd.); Estimated marginal means for interaction / simple-effects follow-up (Lenth RV (2024). emmeans: Estimated Marginal Means, aka Least-Squares Means. R package.); $\eta$\textsuperscript{2} effect-size benchmarks (Cohen, J. (1988). Statistical Power Analysis for the Behavioral Sciences (2nd ed.). Routledge.).

\end{document}